\section{The Finite Difference Method}

\begin{frame}{Explicit Finite Difference Method}
    We use the finite difference method (FDM) to obtain numerical solutions for the valuation of derivatives by solving their underlying differential equations. \newline
    We will use the FDM to obtain numerical solutions to the Black-Scholes equation.
\end{frame}

\begin{frame}{Explicit Finite Difference Method}
    The main goal of the FDM is to convert a continuous differential equation into a set of difference algebraic equations, which can then be solved iteratively.
    \newline
    We will need to define our solution to the problem on the $(S,t)$ plane
\end{frame}

\begin{frame}{Equation Discretization}
    We discretize the equation with respect to $S$ (stock price) and $t$ (time)
    \newline
    Since we know the option to mature at time $T$, we can simply divide $T$ into N equally spaced intervals of $\Delta{t}$ giving $N+1$ points
    $\{0,\Delta{t},2\Delta{t}, ... , N\Delta{t}=T\}$
\end{frame}

\begin{frame}{Equation Discreteization}
    We similarly need to discretize $S$. To do this, we assume that the highest price of the stock is $S_{max}$, where the put option becomes worthless, and since the price can never go below 0, we can split this range into $M$ equally spaced intervals of $\Delta S$.
    \newline
    Since we need the stock price to be large enough, we assume $S_{max}=2S_{0}$, where $S_{0}$ is the initial stock price. Hence we get the set of points,
    $\{0,\Delta{S},2\Delta{S}, ... , M\Delta{S}=S_{max}\}$
\end{frame}


\begin{frame}{Equation Discreteisation - the $(S,t)$ plane}
    \begin{center}
\resizebox{0.65\textwidth}{!}{%
\begin{tikzpicture}[x=1.25cm,y=1cm,>=stealth, every node/.style={font=\normalsize}]

% Axes
\draw[->, thick, QuantNavy] (0,0) -- (6.9,0) node[right, SoftGrey] {Time, \(t\)};
\draw[->, thick, QuantNavy] (0,0) -- (0,6.4) node[above, SoftGrey] {Asset prices, \(S\)};

% Grid points
\foreach \x in {0,1,2,3,4,5,6}{
    \foreach \y in {0,1,2,3,4,5,6}{
        \fill[SoftGrey] (\x,\y) circle (2.2pt);
    }
}

% Axis labels
\node[below, SoftGrey] at (0,0) {\(0\)};
\node[below, SoftGrey] at (1,0) {\(\Delta t\)};
\node[below, SoftGrey] at (2,0) {\(2\Delta t\)};
\node[below, SoftGrey] at (6,0) {\(T\)};

\node[left, SoftGrey] at (0,1) {\(\Delta S\)};
\node[left, SoftGrey] at (0,2) {\(2\Delta S\)};
\node[left, SoftGrey] at (0,6) {\(S_{\max}\)};

\end{tikzpicture}
}
\end{center}

\end{frame}